\subsection{Результати ручного оцінювання якості перекладу}

Для детального аналізу якості машинного перекладу OPUS-MT (baseline-модель)
проведено ручну анотацію помилок із класифікацією за типами~\cite{mqm}:
пропуски (omissions), спотворення (distortions) та додавання (additions). Додатково оцінено
міжфразову когерентність між послідовними сегментами перекладу. Анотація
виконана одним анотатором, що обмежує надійність отриманого результату
(детальніше в розділі <<Висновки>>).

\subsubsection{Класифікація помилок перекладу}

За результатами порівняння машинного перекладу (MT) з референсним людським
перекладом ідентифіковано 114 помилок різних типів. Розподіл помилок за
категоріями наведено в таблиці~\ref{tab:error_distribution}.

\begin{table}[H]
\centering
\caption{Розподіл помилок перекладу за типами}
\label{tab:error_distribution}
\begin{tabular}{|l|c|c|c|}
\hline
\textbf{Тип помилки} & \textbf{Кількість} & \textbf{\%} & \textbf{$R$-коефіцієнт} \\
\hline
Спотворення (distortions) & 91 & 79,8\% & $R_{\text{dist}} = 1,00$ \\
Пропуски (omissions) & 19 & 16,7\% & $R_{\text{omit}} = 0,83$ \\
Додавання (additions) & 4 & 3,5\% & $R_{\text{add}} = 0,33$ \\
\hline
\textbf{Загалом} & \textbf{114} & \textbf{100\%} & --- \\
\hline
\end{tabular}
\end{table}

Коефіцієнт $R$ (rate) обчислено як частку сегментів, у яких виявлено
щонайменше одну помилку відповідного типу: $R = n_{\text{affected}} / N$,
де $N = 12$ --- загальна кількість сегментів. Значення $R_{\text{dist}} = 1,00$
свідчить про наявність спотворень у всіх 12 сегментах корпусу. Високе значення
$R_{\text{omit}} = 0,83$ вказує на систематичні пропуски інформації в більшості
сегментів.

\subsubsection{Категорії спотворень}

Детальний аналіз 91 спотворення дозволив виокремити шість основних категорій
помилок (таблиця~\ref{tab:distortion_categories}).

\begin{table}[H]
\centering
\caption{Категорії спотворень у машинному перекладі}
\label{tab:distortion_categories}
\begin{tabular}{|l|c|c|p{5cm}|}
\hline
\textbf{Категорія} & \textbf{Кількість} & \textbf{\%} & \textbf{Приклад} \\
\hline
Власні назви & 42 & 46,2\% & Rustema Mirov $\rightarrow$ Rustem Umerov \\
Географічні назви & 15 & 16,5\% & Crete $\rightarrow$ Kryvyi Rih \\
Семантичні помилки & 14 & 15,4\% & our killers $\rightarrow$ our warriors \\
Військова термінологія & 12 & 13,2\% & 425th ice storm regiment $\rightarrow$ 425th Assault Battalion \\
Абревіатури & 7 & 7,7\% & MBI $\rightarrow$ Ministry of Internal Affairs \\
Одиниці виміру & 1 & 1,1\% & 20,000 dollars $\rightarrow$ 20,000 hryvnias \\
\hline
\end{tabular}
\end{table}

\paragraph{Власні назви (46,2\%).}
Найбільшу частку спотворень становлять помилки у власних назвах. Модель ASR
некоректно розпізнає українські імена та прізвища, що призводить до каскадного
ефекту в MT. Характерні приклади:
\begin{itemize}[nosep]
    \item <<Rustem Umerov>> $\rightarrow$ <<Rustema Mirov>>;
    \item <<Yevhenii Khmara>> $\rightarrow$ <<Eugene Cloud>>;
    \item <<Keir Starmer>> $\rightarrow$ <<Cyrus Starmer>>;
    \item <<Andrii Sybiha>> $\rightarrow$ <<Andriy Sebiges>>.
\end{itemize}

\paragraph{Географічні назви (16,5\%).}
Значну проблему становить переклад українських топонімів. Виявлено критичні
помилки, коли назви українських міст замінювалися на географічно віддалені
об'єкти:
\begin{itemize}[nosep]
    \item <<Kryvyi Rih>> $\rightarrow$ <<Crete>> (грецький острів);
    \item <<Kryvyi Rih>> $\rightarrow$ <<Crimea>> (політично чутлива територія);
    \item <<Mykolaiv region>> $\rightarrow$ <<Nicholasland>>;
    \item <<Poltava region>> $\rightarrow$ <<Valvas region>>.
\end{itemize}

\paragraph{Семантичні помилки (15,4\%).}
До цієї категорії належать спотворення, що суттєво змінюють зміст повідомлення:
\begin{itemize}[nosep]
    \item <<our warriors>> $\rightarrow$ <<our killers>> (зміна конотації);
    \item <<We also reviewed the situation>> $\rightarrow$ <<The eye doctor figured out>>;
    \item <<resilience>> $\rightarrow$ <<sewers>>;
    \item <<debris clearing>> $\rightarrow$ <<invasion of the wrecks>>.
\end{itemize}

\paragraph{Військова термінологія (13,2\%).}
Помилки в назвах військових підрозділів є критичними для точності повідомлення:
\begin{itemize}[nosep]
    \item <<425th Separate Assault Battalion>> $\rightarrow$ <<425th ice storm regiment>>;
    \item <<414th Unmanned Systems Brigade>> $\rightarrow$ <<414 Tachimadora>>;
    \item <<93rd Separate Mechanized Brigade>> $\rightarrow$ <<93 individual mechanized team>>.
\end{itemize}

\subsubsection{Критичні помилки}

Серед виявлених спотворень ідентифіковано 7 критичних помилок, які суттєво
впливають на сприйняття повідомлення або містять фактичні неточності
(таблиця~\ref{tab:critical_errors}).

\begin{table}[H]
\centering
\caption{Критичні помилки машинного перекладу}
\label{tab:critical_errors}
\small
\begin{tabular}{|p{2.5cm}|p{4cm}|p{4cm}|p{3cm}|}
\hline
\textbf{Сегмент} & \textbf{MT} & \textbf{Референс} & \textbf{Вплив} \\
\hline
В\_Еміратах seg002 & 20,000 dollars & 20,000 hryvnias & Фактична помилка (валюта) \\
\hline
Діалог seg001 & our killers today & our warriors today & Зміна конотації \\
\hline
Знаємо seg001 & Crete & Kryvyi Rih & Географічна плутанина \\
\hline
Наша\_страт. seg001 & Crimea & Kryvyi Rih & Політично чутлива помилка \\
\hline
Отримуємо seg001 & Sewers & resilience & Повне спотворення \\
\hline
В\_Еміратах seg001 & The eye doctor figured out & We also reviewed & Семантична деструкція \\
\hline
Діалог seg002 & 425th ice storm regiment & 425th Assault Battalion & Ідентифікація підрозділу \\
\hline
\end{tabular}
\end{table}

\subsubsection{Аналіз пропусків}

Виявлено 19 пропусків інформації в машинному перекладі. Найчастіше пропускаються:
\begin{itemize}[nosep]
    \item завершальна фраза <<Glory to Ukraine!>> (8 випадків);
    \item деталі військових підрозділів та операцій (5 випадків);
    \item уточнюючі коментарі та оцінки (4 випадки);
    \item назви конкретних організацій (2 випадки).
\end{itemize}

Систематичний пропуск завершальної фрази пов'язаний з особливостями сегментації
та можливим обрізанням кінця аудіофрагменту під час розпізнавання.

\subsubsection{Оцінка міжфразової когерентності}

Для оцінки міжфразової когерентності проаналізовано 5 пар послідовних
сегментів (seg001--seg002) за 5-бальною шкалою: 1~--- відсутність
зв'язності між сегментами, 2~--- суттєві порушення (критичні помилки
в назвах/термінах), 3~--- помірна зв'язність (тематика збережена, але
локальні неточності), 4~--- висока зв'язність (незначні неточності),
5~--- повна когерентність. Результати наведено в таблиці~\ref{tab:coherence}.

\begin{table}[H]
\centering
\caption{Оцінка міжфразової когерентності}
\label{tab:coherence}
\begin{tabular}{|l|c|p{6cm}|}
\hline
\textbf{Пара сегментів} & \textbf{Оцінка} & \textbf{Коментар} \\
\hline
В\_Еміратах & 3 & Тематична послідовність збережена, але непослідовність імен знижує зв'язність \\
\hline
Діалог\_з\_Америкою & 2 & Суттєві термінологічні зміни, критичні помилки в назвах підрозділів \\
\hline
Знаємо\_що\_росіяни & 3 & Спотворення географічних назв впливають на просторову когерентність \\
\hline
Наша\_стратегія & 4 & Краща когерентність, менше критичних помилок \\
\hline
Поведінка & 4 & Відносно висока когерентність, незначні термінологічні неточності \\
\hline
\textbf{Середнє} & \textbf{3,2} & --- \\
\hline
\end{tabular}
\end{table}

Середній показник когерентності 3,2 з 5 вказує на помірну зв'язність між
сегментами. Основними факторами, що знижують когерентність, є:
\begin{itemize}[nosep]
    \item непослідовність у передачі власних назв (одна особа може мати різні
    варіанти імені в різних сегментах);
    \item спотворення географічних назв, що ускладнює просторову орієнтацію;
    \item термінологічні варіації при позначенні тих самих військових підрозділів.
\end{itemize}

\subsubsection{Підсумок ручного оцінювання}

Ручна анотація виявила такі закономірності:

\begin{enumerate}
    \item \textbf{Домінування спотворень}: 79,8\% помилок --- спотворення,
    переважно у власних назвах (46,2\%) та географічних назвах (16,5\%).

    \item \textbf{Каскадний ефект ASR$\rightarrow$MT}: помилки розпізнавання
    власних назв (Whisper) безпосередньо призводять до некоректного перекладу
    (OPUS-MT), оскільки модель MT не має механізму корекції вхідних помилок.

    \item \textbf{Критичні семантичні помилки}: виявлено 7 помилок, що суттєво
    змінюють зміст повідомлення.

    \item \textbf{Помірна когерентність}: середня оцінка 3,2 з 5 свідчить про
    збереження загальної тематичної структури при наявності локальних порушень
    зв'язності.

    \item \textbf{Проблемні категорії}: найскладнішими для системи є власні
    назви, географічні назви та військова термінологія --- категорії, недостатньо
    представлені в навчальних даних моделей.
\end{enumerate}

